\documentclass[12pt]{article}
\usepackage[utf8]{inputenc}
\usepackage[T1]{fontenc}
\usepackage{hyperref}
\title{Maya Samo Language Resources for K-12 Curriculum Development}
\author{Compiled by Claude Sonnet 5.5 (claude-sonnet-5-5) for Andr\'as Kornai}
\date{}
\begin{document}
\maketitle

\section{Introduction}
Maya Samo (ISO 639-3: sym; Glottocode maya1281) is a Mande language of northwestern Burkina Faso, spoken mainly in Sourou Province and neighbouring provinces. Sources differ on classification and on the number of speakers: one search-result summary of the Global Recordings pages gives 38,000, while the reference description for this survey gives about 100,000. Little dedicated educational material was found in web search. This dossier lists only what the search returned.

\section{Primary Educational Resources}
No primers, textbooks or school programmes for Maya Samo were found in search results. The language is not known to be taught in schools.

\section{Linguistic and Reference Materials}
\begin{itemize}
\item Wikipedia, ``Samo language (Burkina)'': general overview of the Samo dialect cluster, noting low intelligibility between varieties. \url{https://en.wikipedia.org/wiki/Samo_language_(Burkina)}
\end{itemize}
No dictionary or grammar of Maya Samo was found in the search results.

\section{Religious and Community Resources}
\begin{itemize}
\item Scripture Earth page for sym, listing scripture resources including audio on Bible.is. \url{https://www.scriptureearth.org/00eng.php?iso=sym}
\item Joshua Project profile of the Maya Samo people in Burkina Faso. \url{https://joshuaproject.net/people_groups/14716}
\end{itemize}

\section{Digital and Audio Resources}
\begin{itemize}
\item Global Recordings Network, Maya Samo (Gomboro variety). Its page says no recordings are currently available. \url{https://www.globalrecordings.net/en/language/16269}
\item Global Recordings Network, Maya Samo (Bangassogo variety). \url{https://www.globalrecordings.net/en/language/16267}
\item Global Recordings Network, Maya Samo (Bounou variety). \url{https://globalrecordings.net/en/language/16268}
\end{itemize}

\section{Curriculum Development Context}
The language is vigorous among its community but has almost no published educational or reference material that could be found. A K-12 curriculum would need new work on orthography, literacy materials and teacher training, and would need to take into account the dialectal diversity of the Samo varieties. Any such work should involve the speaker community directly.

\end{document}